\scanpage{111}

Si $n=2$, ceci donne

\[\left\{\begin{aligned}\varepsilon&=\partial(-1)\\ \unsure{\eta_0|y_0}&=\partial(\widetilde\rho)=\partial(\underbrace{4-\beta_0-\beta_1}_{-(\alpha_0+\alpha_1)})=\partial[-(\alpha_0+\alpha_1)]\\ \eta_1&=\partial\widetilde\sigma=\partial(4-\beta_1)=\partial(2-\alpha_1)\\ \eta_2&=\partial\beta_1=\partial(2+\alpha_1)\end{aligned}\right.\]

\[w_* f_V=(1+\eta_0+\eta_1)(1+\eta_1+\eta_2)\]

avec hyp. parasite que $\widetilde\sigma\deleted{\unsure{\ne}}=2-\alpha_1$ \deleted{\unsure{et $\beta_1=2+\alpha_1$}} \inserted{soit} inversible. Un calcul plus satisfaisant, sans cette hypothèse, s'obtient en prenant l'orthogonal de $\unsure{\underline O_S|\mathcal O_S}(a_0,a_2)$ engendré par $u=\beta_0 a_0+2a_1+a_2$, on trouve $f(u)=-\beta_0(\alpha_0+\alpha_1)$, et $f(a_0)=1$, $f(a_2)=\beta_0\beta_1$ syst. \unsure{orthonormal|orthonormé|orthogonal}. On pose alors

\[\begin{aligned}\varepsilon&=\partial(-1)\\\eta_0&=\partial(\beta_0)\\\eta_1&=\partial(\beta_1)\\\eta_2&=\partial[-(\alpha_0+\alpha_1)]\end{aligned}\]

et on trouve

\[\begin{aligned}w_*(f_V)&=(1+\eta_0+\eta_2)(1+\eta_0\deleted{\unsure{)}}+\eta_1)\qquad\text{i.e.}\\w_1(f_V)&=\deleted{\unsure{\eta_0+}}\eta_1+\eta_2\\w_2(f_V)&=\deleted{\unsure{\eta_1^2}}(\eta_0+\eta_1)(\eta_0+\eta_2)=\eta_0^2+(\eta_0\eta_1+\eta_1\eta_2+\eta_2\eta_0)\end{aligned}\]

% archive sheet number: 110
